\section{System Model and Problem Formulation}

\subsection{Network Model}

We consider an IoT edge network with a set of edge nodes \(V\). In the WLAN setting, each node corresponds to an AP-like edge cache. The mobility-derived topology is represented by a weighted graph
\[
G=(V,E,W),
\]
where \(E\) is the set of cooperative relationships and \(W=\{w_{ij}\}\) is the interaction weight matrix. Larger \(w_{ij}\) indicates stronger mobility-induced relation between nodes \(i\) and \(j\), such as frequent handoffs.

The system contains a content catalog \(\mathcal{C}\). Each node \(i\) has cache capacity \(K_i\), measured as number of objects in this implementation. The binary placement variable is
\[
x_{i,c} =
\begin{cases}
1, & \text{if node } i \text{ stores content } c,\\
0, & \text{otherwise.}
\end{cases}
\]
The capacity constraint is
\[
\sum_{c \in \mathcal{C}} x_{i,c} \le K_i, \qquad \forall i \in V.
\]

\subsection{Perturbed Topology}

Real edge networks are not perfectly stable. We model uncertainty by sampling a perturbed graph
\[
\tilde{G}=(\tilde{V},\tilde{E},\tilde{W}) \sim \mathcal{P}(G).
\]
In the implementation, nodes are removed with probability \(p\) and edges are removed with probability \(p/2\). The largest connected component is retained for evaluation. This models temporary AP unavailability and loss of cooperative links.

\subsection{Access Cost}

Let \(P_{i,c}\) denote the probability that node \(i\) requests content \(c\). Let \(L_{\text{local}}\) be local-cache latency and \(L_{\text{origin}}\) be origin-server latency. Let \(L(i,j)\) be the latency of retrieving content from node \(j\) to node \(i\), computed as a non-decreasing function of shortest-path distance in \(\tilde{G}\). The best non-local retrieval latency is
\[
A_{i,c}(x,\tilde{G})
=
\min\left(
L_{\text{origin}},
\min_{\substack{j \in \tilde{V}\\j \ne i\\x_{j,c}=1}} L(i,j)
\right).
\]
The expected access cost is
\[
C_{\text{access}}(x,\tilde{G})
=
\sum_{i \in \tilde{V}}\sum_{c \in \mathcal{C}}
P_{i,c}
\left[
x_{i,c}L_{\text{local}}+
(1-x_{i,c})A_{i,c}(x,\tilde{G})
\right].
\]

\subsection{Replication, Redundancy, and Relocation}

Replication cost captures the overhead of storing objects:
\[
C_{\text{rep}}(x,\tilde{G})
=
\sum_{i \in \tilde{V}}\sum_{c \in \mathcal{C}}
\gamma f_c x_{i,c}.
\]
Topology-aware redundancy penalizes storing the same object at strongly related nodes:
\[
C_{\text{red}}(x,\tilde{G})
=
\sum_{c \in \mathcal{C}}\sum_{\substack{i,j \in \tilde{V}\\i<j}}
\tilde{w}_{ij}x_{i,c}x_{j,c}.
\]
Relocation cost measures how much a placement differs from a reference placement \(x^0\):
\[
C_{\text{rel}}(x,x^0)
=
\sum_{i \in \tilde{V}}\sum_{c \in \mathcal{C}}
\delta_c |x_{i,c}-x^0_{i,c}|.
\]

\subsection{Optimization Objective}

The stochastic topology-aware cooperative caching problem is
\[
\min_x
\mathbb{E}_{\tilde{G}\sim\mathcal{P}(G)}
\left[
C_{\text{access}}
+ \lambda C_{\text{rep}}
+ \beta C_{\text{red}}
+ \mu C_{\text{rel}}
\right],
\]
subject to binary placement and cache-capacity constraints. The objective captures the central trade-off: aggressive replication can reduce origin misses but wastes storage and increases redundancy; conservative placement reduces storage overhead but may increase access cost; frequent relocation adapts to topology changes but may be operationally expensive.

